% Строчки для суммарной таблицы, в следующей строке прописаны используемые экземпляры узлов
%===s control, LLN0, PDIF1, PDIF2, H2FPHAR1, H5FPHAR1, BLKRCTR1, RCTR1, RMXU1, RMXU2, RMXU3 
\begin{flushleft}
\needspace{4\baselineskip}
\color{uniblue}{\subsection{\raggedright{Дифференциальная токовая защита трансформатора (ДЗТ)}}}
\color{black}
\nopagebreak
\vspace{3mm}
\end{flushleft}

%%%%%%%%%%%%%%%%%%%%%%%%%%%%%%%%%%%%%%%%%%%%%%%%%%%%%%%%%%%%%%%%%%%%%%  ОБЩИЕ СВЕДЕНИЯ %%%%%%%%%%%%%%%%%%%%%%%%%%%%%%%%%%%%%%%%%%%%%%%%%%

\begin{flushleft}
\color{uniblue}{\subsubsection{Общие сведения}}
\color{black}
\nopagebreak
\end{flushleft}

Функция ДЗТ применяется в качестве основной защиты двухобмоточного или трехобмоточного силового трансформатора. ДЗТ обладает абсолютной селективностью и высоким быстродействием. Работа функции не зависит от режима заземления нейтрали сети, и может быть использована в качестве основной защиты от всех видов коротких замыканий. Функция содержит две независимые друг от друга ступени:
\begin{itemize}
\item дифференциальный токовый орган с торможением (ДТЗт);
\item дифференциальная токовая отсечка (ДТО).
\end{itemize}

%%%%%%%%%%%%%%%%%%%%%%%%%%%%%%%%%%%%%%%%%%%%%%%%%%%%%%%%%%%%%%%%%%%%%%  ПРИНЦИП ДЕЙСТВИЯ %%%%%%%%%%%%%%%%%%%%%%%%%%%%%%%%%%%%%%%%%%%%%%%%%%

\begin{flushleft}
\color{uniblue}{\subsubsection{Принцип действия}}
\color{black}
\nopagebreak
\end{flushleft}

Принцип работы дифференциальной защиты основан на первом правиле Кирхгофа. Защищаемый объект рассматривается как электрический узел, а его связи с электрической сетью как ветви узла.

%%%%%%%%%%%%%%%%%%%%%%%%%%%%%%%%%%%%%%%%%%%%%%%%%%%%%%%%%%%%%%%%%%%%%% ЦИФРОВОЕ ВЫРАВНИВАНИЕ %%%%%%%%%%%%%%%%%%%%%%%%%%%%%%%%%%%%%%%%%%%%%%%%%%

\begin{flushleft}
\color{uniblue}{\subsubsection{Цифровое выравнивание подведенных к устройству токов, исключение утроенного тока нулевой последовательности}}
\color{black}
\nopagebreak
\end{flushleft}

Силовые трансформаторы, как правило, имеют различные группы соединений обмоток. Измерительные трансформаторы тока, установленные по сторонам защищаемого трансформатора, как правило, имеют различные номинальные значения первичных и вторичных токов. Прежде чем выполнить сравнение фазных токов каждой из сторон трансформатора,  необходимо выполнить выравнивание измеряемых токов по модулю (амплитуде), по фазе и (при необходимости) выполнить фильтрацию токов нулевой последовательности.
Устройство по заданным параметрам осуществляет выравнивание измеренных токов сторон осуществляется в соответствии с формулой:

% ФОРМУЛА №1
\begin{equation*}
\begin{bmatrix}
I_{A,n*}\\ 
I_{B,n*}\\ 
I_{C,n*}
\end{bmatrix}= k_{AM,n}\times \begin{bmatrix}
M_{11,n} & M_{12,n} & M_{13,n}\\ 
M_{21,n} & M_{22,n} & M_{23,n}\\ 
M_{31,n} & M_{32,n} & M_{33,n}
\end{bmatrix}\times \begin{bmatrix}
I_{A,n}\\ 
I_{B,n}\\ 
I_{C,n}
\end{bmatrix},
\end{equation*}
\begin{eqexpl}[25mm]
\item{$n$}индекс, обозначающий сторону (плечо) ДЗТ;
\item{$k_{AM,n}$}коэффициент амплитудного выравнивания;
\item{$I_{A,n},I_{B,n},I_{C,n}$}измеренные фазные токи стороны $n$;
\item{$M_{11,n}\dots M_{33,n} $}коэффициенты матрицы компенсации фазового сдвига в соответствие  с группой соединения обмоток трансформатора.
\end{eqexpl}

\vspace{1mm}
\emph{\textls[50]{Выравнивание токов сторон по модулю (амплитуде)}}
\par
\vspace{1mm}
Расчет коэффициента амплитудного выравнивания ($k_{AM,n}$) производится по задаваемым параметрам величин номинальных значений силового трансформатора и измерительных трансформаторов тока в соответствии с выражениями:

% ФОРМУЛА №2
\begin{equation*}
k_{AM,n}= \frac{I_{\textsl{перв,n}}\cdot  I_{\textsl{ном.терм,n}} }{I_{\textsl{баз,n}}\cdot I_{\textsl{втор,n}}\cdot k_{\textsl{сх,n}}},
\end{equation*}
\begin{equation*}
I_{\textsl{баз,n}}=\frac{S_{\textsl{баз}}}{\sqrt{3} \cdot U_{\textsl{баз,n}} },
\end{equation*}
\begin{eqexpl}[25mm]
\item{$n$} индекс, обозначающий сторону (плечо) ДЗТ;
\item{$I_{\textsl{перв,n}}$}номинальный первичный ток трансформатора тока стороны $n$;
\item{$I_{\textsl{втор,n}}$}номинальный вторичный ток трансформатора тока стороны $n$;
\item{$I_{\textsl{ном.терм,n}}$}номинальное значение тока аналоговых входов устройства, подключенных к трансформатору тока стороны $n$;
\item{$k_{\textsl{сх,n}}$}коэффициент, учитывающий схему соединения обмоток трансформатора тока стороны~$n$. Данный параметр принимается равным 1, когда трансформаторы тока стороны $n$ собраны в звезду, или равным $\sqrt{3}$, когда трансформаторы тока стороны $n$ собраны в треугольник. Коэффициенты схемы соединения трансформаторов тока для каждой из сторон задается уставками \textbf{<<KсхемСт1>>}, \textbf{<<KсхемСт2>>}, \textbf{<<KсхемСт3>>};
\item{$S_{\textsl{баз}}$}базисная мощность, равная мощности номинальной мощности силового трансформатора;
\item{$U_{\textsl{баз,n}}$}базисное напряжение, равное номинальному напряжению обмотки трансформатора, соответствующей стороне $n$.
\end{eqexpl}

\vspace{1mm}
\emph{\textls[50]{Выравнивание токов сторон по фазе}}
\par
\vspace{1mm}

Для компенсации фазового сдвига, обусловленного различными группами соединения обмоток трансформатора для каждой из сторон предусмотрены задаваемые параметры (\textbf{<<Nсх1>>}, \textbf{<<Nсх2>>}, \textbf{<<Nсх3>>}). Значения данных параметров (от 0 до 11) задают значения коэффициентов матрицы $M_{11,n}\dots M_{33,n}$ в соответствие с таблицей \ref{dzt:tbl0}.

\vspace{1mm}
\emph{\textls[50]{Компенсация утроенного тока нулевой последовательности}}
\par
\vspace{1mm}

Для обмоток силового трансформатора, соединенных в звезду с заземленной нейтралью, и, если не используется других мер исключения токов нулевой последовательности (соединение вторичных обмоток ТТ в треугольник) всегда рекомендуется предусматривать исключение составляющих нулевой последовательности из фазных токов. Компенсация токов нулевой последовательности сторон задается параметрами \textbf{<<Компенс3I0ст1>>}, \textbf{<<Компенс3I0ст2>>}, \textbf{<<Компенс3I0ст3>>}. В случае, когда параметр компенсации токов нулевой последовательности для стороны равен \textbf{<<0>>} (без компенсации тока $3I0$), то в формуле для выравнивания измеряемых токов сторон применяются коэффициенты матрицы \textbf{M} (см. таблицу \ref{dzt:tbl0}). В случае, когда параметр компенсации токов нулевой последовательности для стороны равен \textbf{<<1>>} (с компенсацией токов $3I0$), то в формуле для выравнивания измеряемых токов сторон применяются коэффициенты матрицы \textbf{M0} (см. таблицу \ref{dzt:tbl0}).

%%%%%%%%%%%%%%%%%%%%%%%%%%%%%%%%%%%%%%%%%%%%%%%%%%%%%%%%%%%%%%%%%%%%%%%%%% ТАБЛИЦА 1 %%%%%%%%%%%%%%%%%%%%%%%%%%%%%%%%%%%%%%%%%%%%%%%%%%%%%%%%%%%%%%%%%%%

\begin{longtable}{|>{\centering\arraybackslash}m{2cm}|>{\centering\arraybackslash}m{4cm}|>{\centering\arraybackslash}m{3cm}|>{\centering\arraybackslash}m{4cm}|>{\centering\arraybackslash}m{3cm}|}
\caption{Компенсация группы соединения \hfill\vspace{-0.5\baselineskip}}\label{dzt:tbl0}\\
\hline

\rowcolor{gray!30}
\multicolumn{1}{|c|}{} & \multicolumn{2}{|c|}{Без компенсации токов 3I0} & \multicolumn{2}{|c|}{С компенсацией токов 3I0} \\
\hline
\rowcolor{gray!30}
Векторная группа & Матрица М & Группы соединений обмоток & Матрица М0 & Группы соединений обмоток \\ 
\hline
\endfirsthead
\caption{\emph{Продолжение\hfill\vspace{-0.5\baselineskip}}} \\ % Отступ обозначения таблицы от таблицы и выравнивание надписи слева
\hline

\rowcolor{gray!30}
\multicolumn{1}{|c|}{} & \multicolumn{2}{|c|}{Без компенсации токов 3I0} & \multicolumn{2}{|c|}{С компенсацией токов 3I0} \\
\hline
\rowcolor{gray!30}
Векторная группа & Матрица М & Группы соединений обмоток & Матрица М0 & Группы соединений обмоток \\ 
\hline
\endhead
\hline
\endfoot
\hline
\endlastfoot

\centering 0 & \centering $$\begin{vmatrix}1&0&0\\0&1&0\\0&0&1\end{vmatrix}$$ & \centering \parbox{1cm}{Yy0 Dd0} & \centering $$\displaystyle\frac{1}{3}\times \begin{vmatrix}2&-1&-1\\-1&2&-1\\-1&-1&2\end{vmatrix}$$ & \centering\arraybackslash \parbox{1cm}{Yyn0 YNy0} \\
\hline

\centering 1 & \centering $$\displaystyle\frac{1}{\sqrt{3}}\times\begin{vmatrix}2&-1&-1\\-1&2&-1\\-1&-1&2\end{vmatrix}$$ & \centering \parbox{1cm}{Yd1 Dy1} & \centering Равна матрице №1 & \centering\arraybackslash \parbox{1cm}{YNd1 Dyn1} \\
\hline

\centering 2 & \centering $$\begin{vmatrix}0&-1&0\\0&0&-1\\-1&0&0\end{vmatrix}$$ & \centering \parbox{1cm}{Yy2 Dd2} & \centering $$\displaystyle\frac{1}{3}\times \begin{vmatrix}1&-2&1\\1&1&-2\\-2&1&1\end{vmatrix}$$ & \centering\arraybackslash \parbox{1cm}{YNy2 Yyn2} \\
\hline

\centering 3 & \centering $$\displaystyle\frac{1}{\sqrt{3}}\times\begin{vmatrix}0&-1&1\\1&0&-1\\-1&1&0\end{vmatrix}$$ & \centering \parbox{1cm}{Yd3 Dy3} & \centering Равна матрице №3 & \centering\arraybackslash \parbox{1cm}{YNd3 Dyn3} \\
\hline

\centering 4 & \centering $$\begin{vmatrix}0&0&1\\1&0&0\\0&1&0\end{vmatrix}$$ & \centering \parbox{1cm}{Yy4 Dd4} & \centering $$\displaystyle\frac{1}{3}\times \begin{vmatrix}-1&-1&2\\2&-1&-1\\-1&2&-1\end{vmatrix}$$ & \centering\arraybackslash \parbox{1cm}{YNy4 Yyn4} \\
\hline

\centering 5 & \centering $$\displaystyle\frac{1}{\sqrt{3}}\times\begin{vmatrix}-1&0&1\\1&-1&0\\0&1&-1\end{vmatrix}$$ & \centering \parbox{1cm}{Yd5 Dy5} & \centering Равна матрице №5 & \centering\arraybackslash \parbox{1cm}{YNd5 Dyn5} \\
\hline

\centering 6 & \centering $$\begin{vmatrix}-1&0&0\\0&-1&0\\0&0&-1\end{vmatrix}$$ & \centering \parbox{1cm}{Yy6 Dd6} & \centering $$\displaystyle\frac{1}{3}\times \begin{vmatrix}-2&1&1\\1&-2&1\\1&1&-2\end{vmatrix}$$ & \centering\arraybackslash \parbox{1cm}{YNy6 Yyn6} \\
\hline

\centering 7 & \centering $$\displaystyle\frac{1}{\sqrt{3}}\times\begin{vmatrix}-1&1&0\\0&-1&1\\1&0&-1\end{vmatrix}$$ & \centering \parbox{1cm}{Yd7 Dy7} & \centering Равна матрице №7 & \centering\arraybackslash \parbox{1cm}{YNd7 Dyn7} \\
\hline

\centering 8 & \centering $$\begin{vmatrix}0&1&0\\0&0&1\\1&0&0\end{vmatrix}$$ & \centering \parbox{1cm}{Yy8 Dd8} & \centering $$\displaystyle\frac{1}{3}\times \begin{vmatrix}-1&2&-1\\-1&-1&2\\2&-1&-1\end{vmatrix}$$ & \centering\arraybackslash \parbox{1cm}{YNy8 Yyn8} \\
\hline

\centering 9 & \centering $$\displaystyle\frac{1}{\sqrt{3}}\times\begin{vmatrix}0&1&-1\\-1&0&1\\1&-1&0\end{vmatrix}$$ & \centering \parbox{1cm}{Yd9 Dy9} & \centering Равна матрице №9 & \centering\arraybackslash \parbox{1cm}{YNd9 Dyn9} \\
\hline

\centering 10 & \centering $$\begin{vmatrix}0&0&-1\\-1&0&0\\0&-1&0\end{vmatrix}$$ & \centering \parbox{1cm}{Yy10 Dd10} & \centering $$\displaystyle\frac{1}{3}\times \begin{vmatrix}1&1&-2\\-2&1&1\\1&-2&1\end{vmatrix}$$ & \centering\arraybackslash \parbox{1cm}{YNy10 Yyn10} \\
\hline

\centering 11 & \centering $$\displaystyle\frac{1}{\sqrt{3}}\times\begin{vmatrix}1&0&-1\\-1&1&0\\0&-1&1\end{vmatrix}$$ & \centering \parbox{1cm}{Yd11 Dy11} & \centering Равна матрице №11 & \centering\arraybackslash \parbox{1cm}{YNd11 Dyn11} \\
\hline

\end{longtable}

%%%%%%%%%%%%%%%%%%%%%%%%%%%%%%%%%%%%%%%%%%%%%%%%%%%%%%%%%%%%%%%%%%%%%%  ДИФОРГАН С ТОРМОЖЕНИЕМ %%%%%%%%%%%%%%%%%%%%%%%%%%%%%%%%%%%%%%%%%%%%%%%%%%

\begin{flushleft}
\needspace{3\baselineskip}
\color{uniblue}{\subsubsection{Дифференциальный токовый орган с торможением (ДТЗт)}\label{sec:dtzt}}
\color{black}
\nopagebreak
\end{flushleft}

\par
Характеристика срабатывания \textbf{<<ДТЗт>>} в координатах $I_{\textsl{диф}}$ и $I_{\textsl{торм}}$ представляет собой ломаную линию (см. рисунок \ref{dzt:char}), состоящую из трех участков. Данная характеристика срабатывания обеспечивает чувствительность \textbf{<<ДТЗт>>} при малых токах КЗ в зоне действия защиты.

\vspace{1mm}
\begin{figure}[H]
\centering
\includegraphics[width=0.6\textwidth,height=0.6\textheight,keepaspectratio]{\fbpath/17. ДЗТ 35/_latex/img/chardtz.pdf}
\captionof{figure}{Характеристика срабатывания ДЗТ}\label{dzt:char}
\end{figure}

\par
Первый участок выполнен горизонтальным и определяется уставкой минимального дифференциального тока срабатывания \textbf{<<Iср (ДТЗт)>>}. Второй участок представлен в виде наклонной линии и определяется уставками начального тормозного тока \textbf{<<Iт1>>} и коэффициентом торможения\footnote{Коэффициент торможения это отношение приращения дифференциального тока срабатывания к приращению тормозного тока.} \textbf{<<Kт1>>}. Третий участок представлен в виде наклонной линии и определяется уставками начального тормозного тока \textbf{<<Iт2>>} и коэффициентом торможения \textbf{<<Kт2>>}. Уставки для всех трех фаз задаются одинаковыми. Для исключения возможности излишнего срабатывания при насыщении измерительных ТТ на значениях тормозного тока выше \textbf{<<Iт2>>} в \textbf{<<ДТЗт>>} предусмотрена блокировка при внешних КЗ.
\par
Для каждой из фаз непрерывно выполняется расчет дифференциального и тормозного токов.

Основной рабочей величиной ДЗТ является дифференциальный ток. В общем виде дифференциальный ток равен модулю геометрической суммы токов основных гармоник плеч, определенных с учетом коэффициентов выравниваний и компенсации токов нулевой последовательности (в случае необходимости) по формуле:

% ФОРМУЛА №3
\begin{equation*}
%\label{eq:formula3}
I_{\textsl{диф.ph}} = \left | \sum_{n=1}^{m}I_{ph,n*} \right |,
\end{equation*}
\begin{eqexpl}[25mm]
\item{$ph$}контролируемая фаза (A,B,C);
\item{$I_{ph,n*}$}вектор тока основной гармоники $n$-го плеча, определенного с учетом коэффициентов выравниваний;
\item{$m$}число сторон (до трёх).
\end{eqexpl}

Дифференциальные токи отдельных фаз определяются (в случае, если используется три стороны защищаемого объекта) следующим образом по выражению:  
\begin{equation*}
%\label{eq:formula4}
\begin{split}
I_{\textsl{диф.A}} = \left | I_{A,1*}+I_{A,2*}+I_{A,3*} \right |,\\
I_{\textsl{диф.B}} = \left | I_{B,1*}+I_{B,2*}+I_{B,3*} \right |,\\
I_{\textsl{диф.C}} = \left | I_{C,1*}+I_{C,2*}+I_{C,3*} \right |,
\end{split}
\end{equation*}

В общем виде тормозной ток \textbf{<<ДТЗт>>} равен арифметической полусумме входных токов плеч, определенных с учетом коэффициентов выравниваний по формуле:
\begin{equation*}
%\label{eq:formula5}
I_{\textsl{торм.ph}} = \frac{1}{2} \sum_{n=1}^{m}\left |I_{ph,n*} \right |,
\end{equation*}
\begin{eqexpl}[25mm]
\item{$ph$}контролируемая фаза (A,B,C);
\item{$I_{ph,n*}$}вектор тока основной гармоники n-го плеча, определенного с учетом коэффициентов выравниваний;
\item{$m$}число сторон (до трёх).
\end{eqexpl}

Тормозные токи отдельных фаз определяются (в случае, если используется три стороны защищаемого объекта) следующим образом по выражению:  
\begin{equation*}
%\label{eq:formula4}
\begin{split}
I_{\textsl{торм.A}} = \frac{1}{2}(\left | I_{A,1*}+I_{A,2*}+I_{A,3*} \right |),\\
I_{\textsl{торм.B}} = \frac{1}{2}(\left | I_{B,1*}+I_{B,2*}+I_{B,3*} \right |),\\
I_{\textsl{торм.C}} = \frac{1}{2}(\left | I_{C,1*}+I_{C,2*}+I_{C,3*} \right |),
\end{split}
\end{equation*}

В случае попадания расчетных величин в зону срабатывания на время большее времени уставки формируется сигнал срабатывания \textbf{<<ДТЗт>>}.

\par
Ввод \textbf{<<ДТЗт>>} в работу на этапе параметрирования устройства осуществляется программным переключателем \textbf{SGF1 <<Ввод\_функции>>}. Функция \textbf{<<ДТЗт>>} может быть выведена из работы оперативно путем активации сигнала \textbf{<<ОВ ДТЗт>>}, для перевода действия на сигнал необходимо активировать сигнал \textbf{<<ДТЗт~на~сигн>>} (например, от функциональной клавиши).

\par
Программным переключателем \textbf{SGF2 <<Реж\_блок>>} возможен выбор режима блокировки \textbf{<<ДТЗт>>}: \textbf{<<Без блокировки>>}, \textbf{<<Блокировка по второй гармонике>>}, \textbf{<<Блокировка по пятой гармонике>>}, \textbf{<<Блокировка по второй и пятой гармоникам>>}.

\par
В алгоритме функции \textbf{<<ДТЗт>>} предусмотрена работа с дополнительным подтверждением от функции блокировки от внешних КЗ (\textbf{<<БВКЗ>>}) путем выбора состояния \textbf{<<С контролем БВКЗ>>} программного переключателя \textbf{SGF3 <<Контр\_БВКЗ>>}, описание работы функции \textbf{<<БВКЗ>>} приведено в подразделе \ref{sec:bvkz}.

\par
При выявлении неисправности в цепях ТТ (срабатывание функции \textbf{<<КЦТ>>} по сторонам трансформатора (см. подраздел \ref{sec:kzt}) автоматически увеличивается уровень тока срабатывания горизонтального участка тормозной характеристики для снижения вероятности ложного срабатывания \textbf{<<ДТЗт>>} в режиме протекания токов нагрузки (происходит загрубление уставки \textbf{<<ДТЗт>>} \textbf{<<Iср (ДТЗт)>>} на \textbf{<<Iср (ДТЗт загр)>>} (см. рисунок \ref{dzt:char}).

\par
Алгоритм работы \textbf{<<ДТЗт>>} выполнен в соответствии с рисунком \ref{dzt:img1}. В таблице \ref{dzt:tbl1} приведены параметры, необходимые для настройки функции \textbf{<<ДТЗт>>}.

\vspace{1mm}
\begin{figure}[H]
\includegraphics[width=\textwidth,height=\textheight,keepaspectratio]{\fbpath/17. ДЗТ 35/_latex/img/dtz.pdf}
\captionof{figure}{Функциональная схема дифференциального токового органа с торможением}
\label{dzt:img1}
\end{figure}

\begin{longtable}{|>{\centering\arraybackslash}m{5.4cm}|>{\centering\arraybackslash}m{3.6cm}|>{\centering\arraybackslash}m{4.2cm}|>{\centering\arraybackslash}m{1.8cm}|>{\centering\arraybackslash}m{1cm}|}
\caption{Параметры для настройки функции \textbf{<<ДТЗт>>}\hfill\vspace{-0.5\baselineskip}}\label{dzt:tbl1}\\ 
\hline
\rowcolor{gray!30}
Параметр (Параметр на ИЧМ) & Условное обозначение на схеме & Значение/ Диапазон & Единица измерения & Шаг \\ 
\hline
\endfirsthead
\caption{\emph{Продолжение\hfill\vspace{-0.5\baselineskip}}} \\ % Отступ обозначения таблицы от таблицы и выравнивание надписи слева
\hline
\rowcolor{gray!30}
Параметр (Параметр на ИЧМ) & Условное обозначение на схеме & Значение/ Диапазон & Единица измерения & Шаг \\ 
\hline
\endhead
\hline
\endfoot
\hline
\endlastfoot
\centering Ввод функции в работу (Ввод\_функции) & \centering SGF1 & \centering 0 = Не предусмотрено \\1 = Предусмотрено & \centering -- & \centering\arraybackslash -- \\ 
\hline
\centering Режим блокировки (Реж\_блок) & \centering SGF2 & \centering 1 = Без блокировки \\2 = Блокировка по 2 гармонике \\3 = Блокировка по 5 гармонике \\4 = Блокировка по 2 и 5 гармоникам & \centering -- & \centering\arraybackslash -- \\ 
\hline
\centering Контроль от БВКЗ (Контр\_БВКЗ) & \centering SGF3 & \centering 0 = Не предусмотрено \\1 = Предусмотрено & \centering -- & \centering\arraybackslash -- \\ 
\hline
\centering Начальный дифференциальный ток срабатывания (Iср) & \centering -- & \centering 0,20…0,60 & \centering о.е. & \centering\arraybackslash 0,01 \\ 
\hline
\centering Начальный дифференциальный ток срабатывания в режиме загрубления (Iср.загр) & \centering -- & \centering 0,20…10,00 & \centering о.е. & \centering\arraybackslash 0,01 \\ 
\hline
\centering Коэффициент торможения первого наклонного участка (Kт1) & \centering -- & \centering 0,2…1,0 & \centering -- & \centering\arraybackslash 0,1 \\ 
\hline
\centering Начальный тормозной ток первого наклонного участка (Iт1) & \centering -- & \centering 0,60…3,00 & \centering о.е. & \centering\arraybackslash 0,01 \\ 
\hline
\centering Коэффициент торможения второго наклонного участка (Kт2) & \centering -- & \centering 0,2…1,0 & \centering -- & \centering\arraybackslash 0,1 \\ 
\hline
\centering Начальный тормозной ток второго наклонного участка (Iт2) & \centering -- & \centering 1,20…10,00 & \centering о.е. & \centering\arraybackslash 0,01 \\ 
\hline
\centering Выдержка времени срабатывания (Тср) & \centering Т1 & \centering 0,00…3,00 & \centering с & \centering\arraybackslash 0,01 \\ 
\hline
\end{longtable}

%%%%%%%%%%%%%%%%%%%%%%%%%%%%%%%%%%%%%%%%%%%%%%%%%%%%%%%%%%%%%%%%%%%%%%  ДИФОТСЕЧКА %%%%%%%%%%%%%%%%%%%%%%%%%%%%%%%%%%%%%%%%%%%%%%%%%%
\begin{flushleft}
\needspace{3\baselineskip}
\color{uniblue}{\subsubsection{Дифференциальная токовая отсечка (ДТО)}}
\color{black}
\nopagebreak
\end{flushleft}

\par
Функция \textbf{<<ДТО>>} предназначена для быстрого и селективного отключения повреждений, сопровождающихся большими токами КЗ в зоне действия защиты ДТЗ. Характеристика срабатывания \textbf{<<ДТО>>} (см. рисунок \ref{dzt:char}) представляет из себя прямую линию, определяемую уставкой \textbf{<<Iср (ДТО)>>}. Функция \textbf{<<ДТО>>} работает без каких-либо блокировок и не имеет торможения. Срабатывание Функции \textbf{<<ДТО>>} происходит, если дифференциальный ток превышает уставку \textbf{<<Iср (ДТО)>>}. Возврат Функции \textbf{<<ДТО>>} происходит, если дифференциальный ток ниже уставки \textbf{<<Iср (ДТО)>>} с учетом коэффициента возврата (\textbf{<<Kвоз>>} равен 0,95).

\par
Ввод \textbf{<<ДТО>>} в работу на этапе параметрирования устройства осуществляется программным переключателем \textbf{SGF1 <<Ввод\_функции>>}. Функция \textbf{<<ДТО>>} может быть выведена из работы оперативно путем активации сигнала \textbf{<<ОВ ДТО>>}, для перевода действия на сигнал необходимо активировать сигнал \textbf{<<ДТО~на~сигн>>} (например, от функциональной клавиши).

\par
Алгоритм работы функции \textbf{<<ДТО>>} выполнен в соответствии с рисунком \ref{dzt:img2}. В таблице \ref{dzt:tbl2} приведены параметры, необходимые для настройки функции \textbf{<<ДТО>>}.

\vspace{1mm}
\begin{figure}[H]
\includegraphics[width=\textwidth,height=\textheight,keepaspectratio]{\fbpath/17. ДЗТ 35/_latex/img/dto.pdf}
\captionof{figure}{Функциональная схема \textbf{<<ДТО>>}}
\label{dzt:img2}
\end{figure}

\begin{longtable}{|>{\centering\arraybackslash}m{5.4cm}|>{\centering\arraybackslash}m{3.6cm}|>{\centering\arraybackslash}m{4.2cm}|>{\centering\arraybackslash}m{1.8cm}|>{\centering\arraybackslash}m{1cm}|}
\caption{Параметры для настройки функции \textbf{<<ДТО>>}\hfill\vspace{-0.5\baselineskip}}\label{dzt:tbl2}\\ 
\hline
\rowcolor{gray!30}
Параметр (Параметр на ИЧМ) & Условное обозначение на схеме & Значение/ Диапазон & Единица измерения & Шаг \\ 
\hline
\endfirsthead
\caption{\emph{Продолжение\hfill\vspace{-0.5\baselineskip}}} \\ % Отступ обозначения таблицы от таблицы и выравнивание надписи слева
\hline
\rowcolor{gray!30}
Параметр (Параметр на ИЧМ) & Условное обозначение на схеме & Значение/ Диапазон & Единица измерения & Шаг \\ 
\hline
\endhead
\hline
\endfoot
\hline
\endlastfoot
\centering Ввод функции в работу (Ввод\_функции) & \centering SGF1 & \centering 0 = Не предусмотрено \\1 = Предусмотрено & \centering -- & \centering\arraybackslash -- \\ 
\hline
\centering Дифференциальный ток срабатывания (Iср) & \centering I> & \centering 3,0…20,0 & \centering о.е. & \centering\arraybackslash 0,1 \\ 
\hline
\centering Выдержка времени срабатывания (Тср) & \centering Т1 & \centering 0,00…3,00 & \centering с & \centering\arraybackslash 0,01 \\ 
\hline
\end{longtable}

%%%%%%%%%%%%%%%%%%%%%%%%%%%%%%%%%%%%%%%%%%%%%%%%%%%%%%%%%%%%%%%%%%%%%%  ДЕТЕКТОР 2 ГАРМОНИКИ %%%%%%%%%%%%%%%%%%%%%%%%%%%%%%%%%%%%%%%%%%%%%%%%%%

\begin{flushleft}
\needspace{3\baselineskip}
\color{uniblue}{\subsubsection{Детектор второй гармоники (Д2Г)}}
\color{black}
\nopagebreak
\end{flushleft}

\par
Функция \textbf{<<Д2Г>>} предусмотрена для отстройки ДЗТ от режимов бросков токов намагничивания, обусловленными остаточным насыщением магнитопровода при включении силового трансформатора. Алгоритм функции \textbf{<<Д2Г>>} контролирует уровень второй гармоники в дифференциальном токе и реагирует на отношение модуля второй гармоники дифференциального тока к модулю основной гармоники. Если данное соотношение превышает уставку \textbf{<<Kгарм2/Kгарм1>>} происходит блокировка срабатывания функции \textbf{<<ДТЗт>>}, которая сохраняется до тех пор, пока соотношение не опустится ниже заданной уставки с учетом коэффициента возврата. Срабатывание \textbf{<<ДТЗт>>} блокируется пофазно при появлении соответствующего блокирующего сигнала. Ввод функции \textbf{<<Д2Г>>} в работу осуществляется автоматически при выполнении следующих условий:
\begin{itemize}
\item ввод функции \textbf{<<ДТЗт>>} в работу; 
\item\sloppy посредством программного переключателя \textbf{<<Режим блокировки>>} функции \textbf{<<ДТЗт>>} выбран один из режимов: \textbf{<<Блокировка по 2 гармонике>>} или \textbf{<<Блокировка по 2 и 5 гармоникам>>}.
\end{itemize}
\par
Дополнительно, в логике функции \textbf{<<Д2Г>>} должен быть реализован режим перекрестной блокировки (ввод производится уставкой \textbf{<<Перекр. блок.>>}). Принцип работы данного режима заключается в следующем:
\begin{itemize}
\item при обнаружении условий срабатывания \textbf{<<Д2Г>>} в одной из трех фаз должно произойти блокирование срабатывания \textbf{<<ДТЗт>>} по всем трем фазам на время, задаваемое уставкой \textbf{<<Tперек. бл.>>};
\item при пропадании условий срабатывания \textbf{<<Д2Г>>} от смежных фаз или по истечению времени \textbf{<<Tперек. бл.>>} должно осуществляться блокирование срабатывания \textbf{<<ДТЗт>>} пофазно.
\end{itemize}
\par
Алгоритм работы функции \textbf{<<Д2Г>>} выполнен в соответствии с рисунком \ref{dzt:imgd2g}. В таблице \ref{dzt:tbld2g} приведены параметры, необходимые для настройки функции \textbf{<<Д2Г>>}. 

\vspace{1mm}
\begin{figure}[H]
\includegraphics[width=\textwidth,height=\textheight,keepaspectratio]{\fbpath/17. ДЗТ 35/_latex/img/d2g.pdf}
\captionof{figure}{Функциональная схема \textbf{<<Д2Г>>}}
\label{dzt:imgd2g}
\end{figure}

\begin{longtable}{|>{\centering\arraybackslash}m{5.4cm}|>{\centering\arraybackslash}m{3.6cm}|>{\centering\arraybackslash}m{4.2cm}|>{\centering\arraybackslash}m{1.8cm}|>{\centering\arraybackslash}m{1cm}|}
\caption{Параметры для настройки функции \textbf{<<Д2Г>>}\hfill\vspace{-0.5\baselineskip}}\label{dzt:tbld2g}\\ 
\hline
\rowcolor{gray!30}
Параметр (Параметр на ИЧМ) & Условное обозначение на схеме & Значение/ Диапазон & Единица измерения & Шаг \\ 
\hline
\endfirsthead
\caption{\emph{Продолжение\hfill\vspace{-0.5\baselineskip}}} \\ % Отступ обозначения таблицы от таблицы и выравнивание надписи слева
\hline
\rowcolor{gray!30}
Параметр (Параметр на ИЧМ) & Условное обозначение на схеме & Значение/ Диапазон & Единица измерения & Шаг \\ 
\hline
\endhead
\hline
\endfoot
\hline
\endlastfoot
\centering Режим перекрестной блокировки (Перекр\_бл) & \centering SGF1 & \centering 0 = Не предусмотрено \\1 = Предусмотрено & \centering -- & \centering\arraybackslash -- \\ 
\hline
\centering Коэффициент отношения второй гармоники к первой гармонике (Kгарм2/Kгарм1) & \centering I2г/I1г> & \centering 0,10…0,50 & \centering -- & \centering\arraybackslash 0,01 \\ 
\hline
\centering Выдержка времени на возврат (Tвоз) & \centering Т1 & \centering 0…100 & \centering мс & \centering\arraybackslash 1 \\ 
\hline
\centering Выдержка времени ввода перекрестной блокировки (Tперек.бл) & \centering Т2 & \centering 0,06…4,00 & \centering с & \centering\arraybackslash 0,01 \\ 
\hline
\end{longtable}

%%%%%%%%%%%%%%%%%%%%%%%%%%%%%%%%%%%%%%%%%%%%%%%%%%%%%%%%%%%%%%%%%%%%%%  ДЕТЕКТОР 5 ГАРМОНИКИ %%%%%%%%%%%%%%%%%%%%%%%%%%%%%%%%%%%%%%%%%%%%%%%%%%

\begin{flushleft}
\needspace{3\baselineskip}
\color{uniblue}{\subsubsection{Детектор пятой гармоники (Д5Г)}}
\color{black}
\nopagebreak
\end{flushleft}

\par
Функция \textbf{<<Д5Г>>} предусмотрена для отстройки ДЗТ от режимов бросков токов намагничивания, обусловленными остаточным насыщением магнитопровода при включении силового трансформатора. Алгоритм функции \textbf{<<Д5Г>>} контролирует уровень второй гармоники в дифференциальном токе и реагирует на отношение модуля второй гармоники дифференциального тока к модулю основной гармоники. Если данное соотношение превышает уставку \textbf{<<Kгарм5/Kгарм1>>} происходит блокировка срабатывания функции \textbf{<<ДТЗт>>}, которая сохраняется до тех пор, пока соотношение не опустится ниже заданной уставки с учетом коэффициента возврата. Срабатывание \textbf{<<ДТЗт>>} блокируется пофазно при появлении соответствующего блокирующего сигнала. Ввод функции \textbf{<<Д5Г>>} в работу осуществляется автоматически при выполнении следующих условий:
\begin{itemize}
\item ввод функции \textbf{<<ДТЗт>>} в работу; 
\item\sloppy посредством программного переключателя \textbf{<<Режим блокировки>>} функции \textbf{<<ДТЗт>>} выбран один из режимов: \textbf{<<Блокировка по 5 гармонике>>} или \textbf{<<Блокировка по 2 и 5 гармоникам>>}.
\end{itemize}
\par
Дополнительно, в логике функции \textbf{<<Д5Г>>} должен быть реализован режим перекрестной блокировки (ввод производится уставкой \textbf{<<Перекр. блок.>>}). Принцип работы данного режима заключается в следующем:
\begin{itemize}
\item при обнаружении условий срабатывания \textbf{<<Д5Г>>} в одной из трех фаз должно произойти блокирование срабатывания \textbf{<<ДТЗт>>} по всем трем фазам на время, задаваемое уставкой \textbf{<<Tперек. бл.>>};
\item при пропадании условий срабатывания \textbf{<<Д5Г>>} от смежных фаз или по истечению времени \textbf{<<Tперек. бл.>>} должно осуществляться блокирование срабатывания \textbf{<<ДТЗт>>} пофазно.
\end{itemize}

\par
Алгоритм работы \textbf{<<Д5Г>>} выполнен в соответствии с рисунком \ref{dzt:imgd5g}. В таблице \ref{dzt:tbld5g} приведены параметры, необходимые для настройки функции \textbf{<<Д5Г>>}.

\vspace{1mm}
\begin{figure}[H]
\includegraphics[width=\textwidth,height=\textheight,keepaspectratio]{\fbpath/17. ДЗТ 35/_latex/img/d5g.pdf}
\captionof{figure}{Функциональная схема \textbf{<<Д5Г>>}}
\label{dzt:imgd5g}
\end{figure}

\begin{longtable}{|>{\centering\arraybackslash}m{5.4cm}|>{\centering\arraybackslash}m{3.6cm}|>{\centering\arraybackslash}m{4.2cm}|>{\centering\arraybackslash}m{1.8cm}|>{\centering\arraybackslash}m{1cm}|}
\caption{Параметры для настройки функции \textbf{<<Д5Г>>}\hfill\vspace{-0.5\baselineskip}}\label{dzt:tbld5g}\\ 
\hline
\rowcolor{gray!30}
Параметр (Параметр на ИЧМ) & Условное обозначение на схеме & Значение/ Диапазон & Единица измерения & Шаг \\ 
\hline
\endfirsthead
\caption{\emph{Продолжение\hfill\vspace{-0.5\baselineskip}}} \\ % Отступ обозначения таблицы от таблицы и выравнивание надписи слева
\hline
\rowcolor{gray!30}
Параметр (Параметр на ИЧМ) & Условное обозначение на схеме & Значение/ Диапазон & Единица измерения & Шаг \\ 
\hline
\endhead
\hline
\endfoot
\hline
\endlastfoot
\centering Режим перекрестной блокировки (Перекр\_бл) & \centering SGF1 & \centering 0 = Не предусмотрено \\1 = Предусмотрено & \centering -- & \centering\arraybackslash -- \\ 
\hline
\centering Коэффициент отношения пятой гармоники к первой гармонике (Kгарм5/Kгарм1) & \centering I5г/I1г> & \centering 0,05…0,40 & \centering -- & \centering\arraybackslash 0,01 \\ 
\hline
\centering Выдержка времени на возврат (Tвоз) & \centering Т1 & \centering 0…100 & \centering мс & \centering\arraybackslash 1 \\ 
\hline
\centering Выдержка времени ввода перекрестной блокировки (Tперек.бл) & \centering Т2 & \centering 0,06…4,00 & \centering с & \centering\arraybackslash 0,01 \\ 
\hline
\end{longtable}

%%%%%%%%%%%%%%%%%%%%%%%%%%%%%%%%%%%%%%%%%%%%%% БЛОКИРОВКА ПРИ ВНЕШНИХ КЗ %%%%%%%%%%%%%%%%%%%%%%%%%%%%%%%%%%%%%%%%%%%%%%%

\begin{flushleft}
\needspace{3\baselineskip}
\color{uniblue}{\subsubsection{Блокировка при внешних КЗ (БВКЗ)}\label{sec:bvkz}}
\color{black}
\nopagebreak
\end{flushleft}

\par
Срабатывание функции \textbf{<<БВКЗ>>} является дополнительным критерием стабильной работы ДЗТ при насыщении ТТ. Общий принцип действия состоит в сравнении фаз токов по сторонам защищаемого объекта. Функция БВКЗ выполняется в пофазном исполнении.
Временные диаграммы работы функции БВКЗ в режимах внешнего и внутреннего КЗ приведены на рисунке \ref{dzt:imgcharbvkz}, для однозначного соответствия точек измерения на функциональной схеме (рис. \ref{dzt:imgbvkz}) и на временных диаграммах (рис. \ref{dzt:imgcharbvkz}) точкам измерения присвоены цифры \textcircled{\raisebox{-0.9pt}{1}}...\textcircled{\raisebox{-0.9pt}{5}}.

На функциональной схеме (рис. \ref{dzt:imgbvkz}) и временных диаграммах (рис. \ref{dzt:imgcharbvkz}) указаны точки измерений 

\emph{\textls[50]{Выбор величин для сравнения}}
\par
Каждый период выполняется вычисление амплитуд токов сторон. Ток стороны с максимальной амплитудой выбирается в качестве первой величины для сравнения ($i1$). В качестве второй величины для сравнения ($i2$) используется сумма мгновенных значений остальных сторон.

\emph{\textls[50]{Сравнение фаз}}
\par
Во время положительной полуволны тока каждой из сравниваемых величин устанавливаются сигналы 1 на определенных элементах схемы. Далее эти сигналы суммируются. При внешних КЗ в сумме сигналов практически отсутствуют паузы (паузы обусловлены пороговыми значениями элементов сравнения). При внутренних КЗ паузы по длительности приближаются к значению времени полупериода (см. рисунок \ref{dzt:imgcharbvkz}). Длительность паузы является критерием блокировки при внешних КЗ. Вторая по счету от начала повреждения длительная пауза дает разрешение на работу ДЗТ.

\vspace{1mm}
\begin{figure}[H]
\centering
\includegraphics[width=1\textwidth,height=1\textheight,keepaspectratio]{\fbpath/17. ДЗТ 35/_latex/img/charbvkz.pdf}
\captionof{figure}{Временные диаграммы работы функции \textbf{<<БВКЗ>>}}
\label{dzt:imgcharbvkz}
\end{figure}

\par
Функция \textbf{<<БВКЗ>>} автоматически вводится в работу при введенной в работу функции \textbf{<<ДТЗт>>} и введенном переключателе \textbf{SGF3 <<Контр\_БВКЗ>>} в составе функции \textbf{<<ДТЗт>>} 

\par
Алгоритм работы \textbf{<<БВКЗ>>} выполнен в соответствии с рисунком \ref{dzt:imgbvkz}. Функция \textbf{<<БВКЗ>>} не требует дополнительной настройки.

\vspace{1mm}
\begin{figure}[H]
\includegraphics[width=\textwidth,height=\textheight,keepaspectratio]{\fbpath/17. ДЗТ 35/_latex/img/bvkz.pdf}
\captionof{figure}{Функциональная схема \textbf{<<БВКЗ>>}}
\label{dzt:imgbvkz}
\end{figure}


%%%%%%%%%%%%%%%%%%%%%%%%%%%%%%%%%%%%%%%%%%%%%%%%%%%%%%%%%%%%%%%%% КОНТРОЛЬ ИСПРАВНОСТИ ТОКОВЫХ ЦЕПЕЙ ПО НЕБАЛАНСУ %%%%%%%%%%%%%%%%%%%%%%%%%%%%%%%%%%%%%%%%%%%%%%%%

\begin{flushleft}
\needspace{3\baselineskip}
\color{uniblue}{\subsubsection{Контроль исправности токовых цепей по величине небаланса (КЦТнеб)}}
\color{black}
\nopagebreak
\end{flushleft}

\par
Функция контроля величины небаланса в дифференциальных цепях является дополнительной. Принцип действия этой функции основан на измерении тока небаланса в дифференциальных цепях: если дифференциальный ток превышает уставку тока небаланса \textbf{<<Iср>>} в течение выдержки времени \textbf{<<Tср>>}, то выдается сигнал о неисправности в сигнализацию.

\par
Алгоритм работы \textbf{<<КЦТнеб>>} выполнен в соответствии с рисунком \ref{dzt:imgkzt}. В таблице \ref{dzt:tblkzt} приведены параметры, необходимые для настройки функции \textbf{<<КЦТнеб>>}.

\vspace{1mm}
\begin{figure}[H]
\includegraphics[width=\textwidth,height=\textheight,keepaspectratio]{\fbpath/17. ДЗТ 35/_latex/img/kzt.pdf}
\captionof{figure}{Функциональная схема \textbf{<<КЦТнеб>>}}
\label{dzt:imgkzt}
\end{figure}

\begin{longtable}{|>{\centering\arraybackslash}m{5.4cm}|>{\centering\arraybackslash}m{3.6cm}|>{\centering\arraybackslash}m{4.2cm}|>{\centering\arraybackslash}m{1.8cm}|>{\centering\arraybackslash}m{1cm}|}
\caption{Параметры для настройки функции \textbf{<<КЦТнеб>>}\hfill\vspace{-0.5\baselineskip}}\label{dzt:tblkzt}\\ 
\hline
\rowcolor{gray!30}
Параметр (Параметр на ИЧМ) & Условное обозначение на схеме & Значение/ Диапазон & Единица измерения & Шаг \\ 
\hline
\endfirsthead
\caption{\emph{Продолжение\hfill\vspace{-0.5\baselineskip}}} \\ % Отступ обозначения таблицы от таблицы и выравнивание надписи слева
\hline
\rowcolor{gray!30}
Параметр (Параметр на ИЧМ) & Условное обозначение на схеме & Значение/ Диапазон & Единица измерения & Шаг \\ 
\hline
\endhead
\hline
\endfoot
\hline
\endlastfoot
\centering Ввод функции в работу (Ввод\_функции) & \centering SGF1 & \centering 0 = Не предусмотрено \\1 = Предусмотрено & \centering - & \centering\arraybackslash - \\ 
\hline
\centering Ток срабатывания контроля небаланса (Iср) & \centering Idiff> & \centering 0,04…2,00 & \centering о.е. & \centering\arraybackslash 0,01 \\ 
\hline
\centering Выдержка времени срабатывания контроля небаланса (Тср) & \centering Т1 & \centering 0,0…110,0 & \centering с & \centering\arraybackslash 0,1 \\ 
\hline
\end{longtable}


